%
% Eindeutigkeit der Erweiterung des Ableitungsoperators
%
\subsection{Eindeutigkeit der Erweiterung des Ableitungsoperators}
Wir betrachten in diesem Abschnitt einen Differentialkörper $F$ mit
der Deriviation $D$ und stellen die Frage, ob sich der Differentialoperator
$D$ auf eindeutige Weise auf eine Erweiterung $F(\vartheta)$ ausdehnen
lässt.

%
% Transzendente Erweiterungen
%
\subsubsection{Transzendente Erweiterungen}
Sei $\vartheta$ eine transzendente Erweiterung von $F$ und
\[
f(\vartheta) = \frac{a(\vartheta)}{b(\vartheta)}
\]
eine rationale Funktion von $\vartheta$ mit Koeffizienten in $F$, dargestellt
als Bruch von Polynomen $a(\vartheta),b(\vartheta)\in F[\vartheta]$.
Wir müssen zeigen, dass es nur eine mögliche Definition einer Deriviation
$\tilde{D}$ auf solchen Brüchen gibt.
$\tilde{D}$ muss alle Rechenregeln einer Deriviation erfüllen, insbesondere
die Quotientenregel
\begin{align*}
\tilde{D}f(\vartheta)
&=
\frac{
\tilde{D}a(\vartheta)\cdot b(\vartheta)-a(\vartheta)\cdot\tilde{D}b(\vartheta)
}{
b(\vartheta)^2
},
\intertext{die Linearität}
\tilde{D} a(\vartheta)
&=
\tilde{D} \sum_{k=1}^n a_k\vartheta^k
=
\sum_{k=1}^n \tilde{D}(a_k\vartheta^k)
\intertext{und die Produktregel}
&=
\sum_{k=1}^n \tilde{D}a_k\cdot \vartheta^k
+
\sum_{k=1}^n a_k\cdot \tilde{D} \vartheta^k
=
\sum_{k=1}^n Da_k\cdot \vartheta^k
+
\sum_{k=1}^n a_k\cdot
k
\vartheta^{k-1}
\tilde{D}
\vartheta.
\end{align*}
Die Derivation $\tilde{D}$ ist also eindeutig bestimmt, sobald
$\tilde{D}\vartheta$ festgelegt ist.

Für eine logarithmische Erweiterung $\vartheta=\log u$ mit $u\in F$ gilt
$\tilde{D}\vartheta=Du/u$.
Da $Du\in F$ und $u\in F$ ist die Ableitung von $\vartheta$ eindeutig
festgelegt und damit $\tilde{D}$ auf ganz $F(\vartheta)$.
%
Für eine exponentielle Erweiterung $\vartheta=e^u$ mit $u\in F$ gilt
$\tilde{D}\vartheta=\tilde{D}e^u = e^uu =\vartheta u$.
Auch in diesem Fall ist die Ableitung von $\vartheta$ eindeutig
festgelegt.
Es folgt, dass es genau eine Erweiterung der Deriviation $D$ auf
$F$ zu einer Derivation $\tilde{D}$ auf ganz $F(\vartheta)$ gibt.

%
% Algebraische Erweiterungen
%
\subsubsection{Algebraische Erweiterungen}
Sei $\vartheta$ eine algebraische Erweiterung mit dem Minimalpolynom
$m(X) \in F[X]$ mit Koeffizienten in $F$.
Auch in diesem Fall stellt sich die Frage, ob die Ausdehnung $\tilde{D}$
auf $F(\vartheta)$ eindeutig bestimmt ist.
Mit zum Fall einer transzendenten Erweiterung analogen Rechnungen folgt
auch hier, dass die Derivation eindeutig festgelegt ist, sobald
$\tilde{D}\vartheta$ festgelegt wird.
Durch Ableiten der Definitionsgleichung
\begin{align}
\sum_{k=0}^n m_k\vartheta^k
&=
0
\intertext{erhält man die Gleichung}
\sum_{k=0}^n Dm_k\vartheta^k
+
\sum_{k=0}^n m_k k\vartheta^{k-1}\tilde{D}\vartheta
&=
0
\notag
\\
\sum_{k=0}^n Dm_k\vartheta^k
+
\biggl(\sum_{k=0}^n m_k k\vartheta^{k-1}\biggr)\tilde{D}\vartheta
&=
0,
\notag
\intertext{die man nach}
\tilde{D}\vartheta
&=
-\frac{\displaystyle
\sum_{k=0}^n Dm_k \cdot \vartheta^k
}{\displaystyle
\sum_{k=0}^n m_kk\vartheta^{k-1}
}
\label{buch:intalg:elementar:Derweiterung:alg}
\end{align}
auflösen kann.
Dies zeigt, dass $\tilde{D}\vartheta$ eindeutig festgelegt ist.

Die zusätzliche Schwierigkeit in diesem Fall ist, dass es mehrere
Funktionen $\vartheta$ gibt, die die Gleichung $m(\vartheta)=0$ erfüllen.
Zum Beispiel definiert das Minimalpolynom $m(X) = X^2-f$ zwei Wurzeln
$\!\sqrt{f}$ und $-\!\sqrt{f}$.
Wir müssen in diesem Fall zusätzlich zeigen, dass
unsere Konstruktion von $\tilde{D}\vartheta$ nicht von der Wahl der
Lösung von $m(\vartheta)=0$ abhängt.
Die Formel~\eqref{buch:intalg:elementar:Derweiterung:alg} zeigt aber,
dass die Ableitung vollständig durch die Koeffizienten des
Minimalpolynoms $m$ festgelegt wird.
Ersetzt man also $\vartheta$ durch eine andere Nullstelle von $m$,
hat die resultierende Ableitung die gleichen algebraischen Eigenschaften.




